\scanpage{67}
\mspage{B.}
% archive foliation: 66
on peut donc supposer que $Z$ est \uline{projectif},
\deleted{Si dim} et $=$ \uline{normal} si on veut.
Si \inserted{$\dim Z\leq0$, trivial. Si}
$\dim Z=1$, on aura $Y=\varnothing$, i.e. $X=Z$,
et le th. est connu. Si $\dim Z\geq2$ on utilise ceci:

\uline{Lemme} de \uline{Bertini} (Zariski)
Soit $k$ corps,
$X\xrightarrow{f}\mathbf P_k^r$ un \uline{morphisme},
avec $X$ \uline{irréductible},
$\dim f(X)=n$, soit $1\leq s\leq n-1$, et soit $u$
le point générique de
$\underbrace{\mathbf P_k^r\times\cdots\times\mathbf P_k^r}_{s}$,
et $K$ une clôture algébrique de $k(u)$,
enfin $L_u$ la variété linéaire de $\mathbf P_K^r$ définie par $u$,
et
\[
X_u=X\times_{\mathbf P_k^r}L_u=X_K\times_{\mathbf P_K^r}L_u.
\]
Alors $X_u$ est irréductible.

\margin{peut-on remplacer ``irréductible'' par ``connexe'' ?
\uline{non}}

\margin{\diagram{\begin{tikzcd}
X'_u\arrow[d]&&X'\arrow[d]\\
X_u\arrow[r]\arrow[d]&X_K\arrow[r]\arrow[d]&X\arrow[d]\\
L_u\arrow[r]&\mathbf P_K^r\arrow[r]&\mathbf P_k^r
\end{tikzcd}}}

\uline{Corollaire} \deleted{\unsure{Soit}}
Si $X$ est \uline{normal}, alors
$\pi_1(X_u)\to\pi_1(X)$ est \uline{surjectif}.

En effet, si $X'$ est un revêtement \deleted{normal}
étale connexe de $X$, étant normal il est \uline{irréductible},
donc $X'_u=X'\times_X X_u$ est irréductible,
donc connexe, OK.

\margin{Si on admet le th. de connexion,
on \deleted{peut} \unsure{pourra alors} utiliser une
\unsure{section}
\deleted{\unsure{pas}}
``\unsure{spéciale}'' \ldots\ pour une courbe.}
